\documentclass{article}
\usepackage{graphicx} % Required for inserting images

\title{Sistemi Operativi}
\author{Tommaso Trentin}
\date{June 2025}

\begin{document}

\newpage
\tableofcontents
\newpage
lololol
\section{CPU Scheduling}
\subsection{Scheduling dei Processi}
\subsubsection{Concetti Generali}
In generale risolvere un problema di scheduling consiste nel \textbf{mappare} un certo numero di \textbf{compiti} su un (diverso) numero di \textbf{esecutori} in modo da realizzare obbiettivi globali del sistema rispettando \textbf{vincoli} prefissati. \\
I principali \textbf{eventi} che possono causare uno \textbf{switch di contesto} (e che fanno quindi entrare in gioco lo scheduler) sono:
\begin{itemize}
  \item interruzione relativa a timer (termine del quanto di tempo)
  \item interruzione dovuta a dispositivi I/O
  \item invocazione di system call (es. exit())
  \item gestione di segnali
\end{itemize}
\subsubsection{Compiti del Dispatcher}
Lo scheduler sceglie il processo da eseguire, successivamente il dispatcher:
\begin{itemize}
    \item effettua il context switch
    \item inizializza program counter alla prossima istruzione del programma da (continuare a) eseguire
    \item effettua il passaggio alla modalità utente
\end{itemize}
Il tempo per fare ciò viene detto latenza di dispatch.
\subsubsection{Parametri Generali}
\begin{itemize}
    \item numero di processori (1+), omogeneità/disomogeneità dei processori (in tipo, velocità, ISA…)
    \item staticità/dinamicità insieme dei task
    \item on-line e off-line scheduling
    \item  tempi di esecuzione noti a priori o sconosciuti
    \item presenza o meno di preemption (prelazione)
    \item presenza o meno di precedenze tra task
    \item presenza o meno di priorità
    \item imposizione o meno di scadenze (deadline) su inizio/terminazione task
    \item presenza o meno di \textbf{obbiettivo} da ottimizzare (tempo di risposta, tempo di completamento, produttività…)
\end{itemize}
\subsubsection{Obbiettivi}
Obbiettivi diversi, spesso antitetici, sono ottimizzazione di:
\begin{itemize}
    \item lunghezza dello schedule
    \item \textbf{tempo di risposta} massimo
    \item \textbf{tempo di completamento} medio
    \item tempo di risposta medio
    \item tempo di attesa
    \item numero minimo di CPU utilizzate
    \item produttività (num. task/tempo $\rightarrow$ \textbf{throughput})
    \item utilizzo della CPU
    \item massimo ritardo rispetto a eventuali deadline
\end{itemize}
\subsubsection{Ctiteri per CPU Scheduling}
Dal punto di vista \textbf{utente}:
\begin{itemize}
    \item minimizzare tempo risposta
    \item minimizzare tempo completamento
    \item rispettare deadline
\end{itemize}
Dal punto di vista \textbf{sistema}:
\begin{itemize}
    \item massimizzare utilizzo CPU
    \item garantire \textbf{fairness} (problema della \textbf{starvation})
    \item massimizzare throughput
    \item riflettere le priorità
\end{itemize}
\textbf{\textit{Sia scheduler che dispatcher rappresentano bottleneck per il S.O.}}
\subsubsection{CPU-burst e I/O-burst}
Il comportamento tipico di un processo consiste nell'alternare sequenze di operazioni di elaborazione (CPU-burst) a sequenze di operazioni di I/O (I/O burst), terminando con un CPU-burst. \\
\includegraphics[width=0.4\textwidth]{OS_images/CPU_IO_Burst.png} \\
Processi CPU-bound avranno pochi CPU-burst lunghi (P1), mentre processi I/O-bound avranno molti CPU-burst brevi (P2): \\
\includegraphics[width=0.6\textwidth]{OS_images/CPU_IO_Bound.png}
\subsubsection{Scheduling Preemptive e Non-preemptive}
Lo schema di scheduling può essere con o senza \textbf{diritto di prelazione}:
\begin{itemize}
    \item \textbf{Nonpreemptive scheduling}: processo running continua a detenere la CPU fino a quando termina o la cede
    \item \textbf{Preemptive scheduling}: processo può essere interrotto (portato a ready) a seguito di un evento “esterno” (scade quanto di tempo, altro processo diventa ready…)
\end{itemize}
Nonpreemptive scheduling tipico dei sistemi batch, preemptive scheduling evita che processi “meno importanti” ostacolino quelli critici (comportando però un overhead di gestione)
\subsection{Algoritmi di Scheduling Processi}
\subsubsection{FCFS - First Come First Served}
\begin{itemize}
    \item Si seleziona il processo che è nella ready-queue da più tempo (processi serviti nell'ordine in cui entrano nella ready-queue)
    \item Nonpreemptive
\end{itemize}
\includegraphics[width=0.6\textwidth]{OS_images/FCFS.png}
\begin{itemize}
    \item Vantaggi
        \begin{itemize}
            \item Semplice da implementare, non comporta overhead di gestione
            \item Non ci può essere starvation in quanto tutti i processi ready, prima o poi, ottengono uso CPU
        \end{itemize}
    \item Svantaggi
        \begin{itemize}
            \item Tempo d’attesa medio molto legato all’ordine di arrivo dei burst
            \item Tempo di risposta di un processo può essere molto alto anche per processi di breve durata
            \item Un processo CPU-bound tende a monopolizzare il processore a discapito dei processi I/O-bound
                \begin{enumerate}
                    \item Processi I/O-bound terminano burst di I/O e si accodano in ready
                    \item CPU-bound rilascia la CPU per fare I/O
                    \item Tutti gli I/O-bound terminano il burst di CPU per fare I/O
                    \item Il CPU-bound termina l’I/O e torna ready (”prima” degli I/O-bound) e tutto si ripete
                \end{enumerate}
            Si tende quindi a sottoutilizzare sia device che CPU (convoy effect)
        \end{itemize}
\end{itemize}
\subsubsection{SJF - Shortest Job First}
\begin{itemize}
    \item Si seleziona il processo che avrà il prossimo CPU-burst più breve
    \item Nonpreemptive
\end{itemize}
\includegraphics[width=0.6\textwidth]{OS_images/SJF.png}
\subsubsection{SRTF - Shortest Remaining Time First}
\begin{itemize}
    \item Si seleziona il processo che avrà il prossimo CPU-burst più breve
    \item Se P1 è running e in ready-queue giunge P2 con durata inferiore al tempo restante di P1, allora P2 sottrae la CPU a P1
    \item Preemptive
\end{itemize}
\includegraphics[width=0.6\textwidth]{OS_images/SRTF.png}
\subsubsection{Caratteristiche SJF / SRTF}
\begin{itemize}
    \item Processi con CPU-burst più breve sopravanzano gli altri
    \item SJF rappresenta lo scheduling ottimale: realizza il tempo di attesa medio migliore
    \item Può esserci starvation: i processi con CPU-burst lunghi tendono a restare molto tempo in ready-queue, non garantito ricevano l’uso della CPU in un tempo limitato
    \item Richiedono di prevedere la durata dei CPU-burst (non sempre possibile determinare queste durate)
\end{itemize}
\subsubsection{Stima durata CPU burst}
La previsione/stima della durata del prossimo CPU-burst di un processo può essere compiuta basandosi sui suoi CPU-burst del passato tramite media esponenziale:
$$\tau_{n+1} = \alpha t_n + (1 - \alpha)\tau_n$$
\\
dove $t_n$ è la durata dell’n-esimo CPU-burst cioè la storia recente, mentre $\tau_n$ rappresenta il contributo della storia remota
Sviluppando il termine n-esimo si ottiene:
$$\tau_{n+1} = \alpha tn + (1 - \alpha)\alpha \tau_{n-1} + (1 - \alpha)^2\alpha \tau_{n-2} + \ldots + (1 - \alpha)^{n+1}\tau_0$$
con $\tau_0$ e $0 \leq \alpha \leq 1$ costanti predefinite. In particolare $\alpha$ permette di modulare il peso della storia recente/remota.
\subsubsection{HRRF - Highest Response Ratio First}
\begin{itemize}
    \item Si seleziona il processo che ha il \textbf{Response Ratio} R maggiore, dove \textbf{R = (ta + ts) / ts} con ta il tempo che il processo ha speso in attesa della CPU e ts la durata prevista del suo CPU-burst
    \item Nonpreemptive
\end{itemize}
\includegraphics[width=0.6\textwidth]{OS_images/HRRF.png}
\begin{itemize}
    \item Processi con CPU-burst brevi sono favoriti ma si tiene conto dell’età del processo, processi con tempo di servizio ts lungo guadagnano priorità man mano che invecchiano
    \item Meccanismo di \textbf{aging}: evita la starvation
    \item Come SJF e SRTF richiede stime su tempo di servizio richiesto da processi
\end{itemize}
\subsubsection{Scheduling con Priorità}
Viene assegnata una priorità ad ogni processo e CPU viene allocata a processo ready con priorità superiore.
Nonpreemptive o preemptive (in questo caso un processo ready con priorità superiore sottrae la CPU ad un processo running di priorità inferiore)
\begin{itemize}
    \item Gestione delle priorità
        \begin{itemize}
            \item Priorità statiche
                \begin{itemize}
                    \item priorità assegnata ad un processo resta costante
                    \item facile implementazione
                    \item minimo overhead di gestione
                    \item no adattamento a cambiamenti del carico di sistema
                \end{itemize}
            \item Priorità dinamiche
                \begin{itemize}
                    \item reattività rispetto a mutazioni del carico di sistema
                    \item permettono di “seguire” evoluzione del comportamento dei processi
                    \item maggior overhead per gestione
                    \item aumento dinamico (solitamente) della priorità ai processi I/O-bound
                \end{itemize}
        \end{itemize}
    \item Priorità e Starvation
        Starvation = processo P, pur essendo ready, non ottiene mai uso della CPU.\\
        In una politica con priorità potrebbe verificarsi se venissero continuamente creati processi con priorità superiore a quella di P.\\
        Per evitare starvation si usa una strategia di aging dei processi (essi tendono a migliorare la loro priorità con il passare del tempo)
\end{itemize}
\subsubsection{Round Robin (circolare)}
Come per FCFS si serve il processo che attende da più tempo \\
Preemptive usando timer: allo scadere del quale il processo running viene rimesso in ready-queue \\
\includegraphics[width=0.6\textwidth]{OS_images/RR.png}
\\
Processi I/O-bound penalizzati (in generale hanno CPU-burst più brevi del quanto di tempo), burst termina prima della fine del quanto (vengono messi in waiting prima di aver sfruttato tutto tempo a disposizione) \\
Processi CPU-bound tendono a esaurire il quanto, quindi anche se mix tra CPU e I/O bound è adeguato in media i CPU-bound sfruttano (tendenzialmente) di più la CPU. Vantaggio: evita la starvation
\subsubsection{Durata quanto di tempo}
Durata di un quanto è critica per il comportamento del sistema:
\begin{itemize}
    \item deve essere sensibilmente superiore al tempo necessario per gestire interruzioni / switch di contesto
    \item si rilevano buoni risultati se durata del quanto è maggiore della durata di almeno 80\% di CPU-burst
    \item non può essere troppo lungo $\rightarrow$ penalizza processi I/O-bound (tende ad un FCFS)
\end{itemize}
\subsubsection{Scheduling con Code Multiple}
Più code di processi:
\begin{itemize}
    \item processi pratizionati in gruppi disgiunti secondo opportuni criteri (es. processi batch e processi interattivi)
    \item processi assegnati permanentemente alle code
    \item ogni coda gestita mediante algoritmo di scheduling più adatto a caratteristiche dei suoi processi
    \item si deve stabilire uno scheduling di code:
        \begin{itemize}
            \item relazione di priorità tra code e strategia preemptive: processo di una coda viene interrotto appena un altro processo entra in una delle code con priorità maggiore
            \item \textbf{time slicing} a livello delle code: ad ogni coda viene assegnata una percentuale di tempo totale di CPU
        \end{itemize}
\end{itemize}
\subsubsection{Scheduling con Code Multiple -Retroazione}
Evoluzione delle code multiple:
\begin{itemize}
    \item \textbf{Evita starvation}, permettendo ai processi di migrare tra code secondo criteri di \textbf{aging}.
    \item Solitamente \textbf{preemptive}.
    \item I processi che utilizzano a lungo la CPU vengono spostati in code con \textbf{priorità più bassa}.
    \item I processi \textbf{interattivi (I/O-bound)} tendono invece a ricevere \textbf{priorità più alta}.
    \item Un processo che resta troppo tempo in una coda a bassa priorità può essere \textbf{promosso} a una coda con priorità maggiore (per evitare starvation).
    \item \textbf{Overhead di gestione} viene bilanciato da una maggiore \textbf{sensibilità e reattività del sistema}.
\end{itemize}
\includegraphics[width=0.6\textwidth]{OS_images/Code_Multiple_4.png}
\\
Parametri da considerare:
\begin{itemize}
    \item Numero code da utilizzare
    \item Algoritmi di scheduling per ogni coda.
    \item \textbf{Regole di migrazione} dei processi tra le code (es. un processo preempted può “retrocedere”, mentre l’\textbf{aging} lo fa salire di priorità).
    \item \textbf{Posizione iniziale} nuovi processi (es. si può scegliere di inserirli direttamente nella coda con priorità più alta).
\end{itemize}
\subsubsection{Comparazione tra algoritmi}
\includegraphics[width=0.8\textwidth]{OS_images/Algorithm_Table.png}
\subsubsection{Fair Share Scheduling}
Processi \textbf{partizionati in gruppi} e scheduling avviene mirando a \textbf{equità di trattamento rispetto ai gruppi}:
\begin{itemize}
    \item adatto a sistemi multi utente
    \item adatto in presenza di applicazioni multi-thread e scheduling viene effettuato sui thread
    \item può trattare diversamente processi di differenti tipologie di utente
\end{itemize}
\subsection{Scheduling dei Thread}
\subsubsection{Kernel-Thread e User-Thread}
\begin{itemize}
    \item User Threads (UT)
        \begin{itemize}
            \item Gestiti nello spazio utente (user space).
            \item \textbf{Il kernel non ne è consapevole}.
            \item Più veloci da creare e gestire.
            \item \textbf{Problema}: se un thread effettua una system call bloccante, blocca tutti i thread del processo
            \item \textbf{PROCESS-CONTENTION SCOPE}
        \end{itemize}
    \item Kernel Threads (KT)
        \begin{itemize}
            \item Gestiti dal kernel, nello spazio kernel.
            \item Ogni thread è schedulato \textbf{indipendentemente}.
            \item Più costosi da gestire in termini di performance.
            \item \textbf{SYSTEM-CONTENTION SCOPE}
        \end{itemize}
\end{itemize}
\subsection{Scheduling per Sistemi Multiprocessore}
\subsubsection{Scheduling in Sistemi Multiprocessore}
In linea di principio possono essere adattate le tecniche di scheduling per i casi single-processor anche in presenza di più unità di elaborazione. \\
Il carico deve però essere distribuito e \textbf{bilanciato}.
\textbf{Tipologie di multielabolarzione}
\begin{itemize}
    \item \textbf{Asimmetrica}: attività di sistema (scheduling, gestione, interruzioni…) affidate tutte a un unico processore \textbf{master server}
        \begin{itemize}
            \item codice degli utenti eseguito dagli altri processori
            \item \textbf{difetto}: master server è un potenziale bottleneck del sistema
        \end{itemize}
    \item \textbf{Simmetrica}: ogni processore ha uno scheduler che attinge alla ready-queue che può essere
        \begin{itemize}
            \item di sistema: possibili race condition se più processori scelgono stesso task (serve mutua esclusione all’interno di ready-queue $\rightarrow$ possibile bottleneck)
            \item per-processor: serve un meccanismo per distribuzione/bilanciamento carico tra ready-queue
        \end{itemize}
\end{itemize}
\textbf{Scheduling (Multi)processori Multicore}
\begin{itemize}
    \item Processori multicore presentano più core in stesso chip fisico
    \item S.O. vede ogni core come un processore logico separato
    \item Ogni core esegue (solitamente) più thread hardware: \textbf{chip multithreading}, \textbf{hyper-threading}
    \item Presenza di multithreading a livello di core comporta doppio livello di scheduling:
        \begin{itemize}
            \item scheduling (hardware) di thread hardware sui core della CPU multicore
            \item scheduling (software) di thread/processi/task sulle CPU (lo scheduling studiato finora)
        \end{itemize}         \includegraphics[width=0.6\textwidth]{OS_images/Multi_Core_Scheduling.png}
\end{itemize}
\subsubsection{Bilanciamento del Carico}
Con più unità di elaborazione va massimizzato distribuendo equamente il carico tra i core del sistema, si utilizzano principalmente due tecniche base (non esclusive):
\begin{itemize}
    \item Push migration: un processo dedicato controlla periodicamente carico di lavoro assegnato a predecessori $\rightarrow$ sottrae task da CPU overloaded e li riassegna a CPU meno cariche
    \item Pull migration: ogni processore, quando inattivo, preleva un task da uno dei processori overloaded
\end{itemize}
\subsubsection{Sistemi Eterogenei}
Presenza di diversi processori con caratteristiche non omogenee:
\begin{itemize}
    \item diverso tipo (processori grafici / dedicati a task specifici)
    \item diverso set di istruzioni
    \item diversa potenza di calcolo / consumo energetico
\end{itemize}
Utile in caso si debba dare maggiore importanza a consumo energetico (sistemi embedded…)
\subsection{Scheduling Real-time}
\subsubsection{Sistemi Real-Time}
\begin{itemize}
    \item Utilizzati in dispositivi di largo consumo (sistemi embedded, player, robotica, domotica)
    \item Funzionalità devono rispettare vincoli temporali (scadenze, durate, tempi di reazione…)
    \item Obbiettivo primario rispetto di deadline e vincoli ai tempi di reazione
    \item Necessitano algoritmi di scheduling specifici
\end{itemize}
\textbf{Tipologie di Sistemi Real-Time}:
\begin{itemize}
    \item Real-time soft: ai task critici viene concessa precedenza (no violazioni deadline però)
    \item Real-time hard: task devono necessariamente rispettare deadline
\end{itemize}
Solitamente sono event-driven (all’accadere di un evento il sistema reagisce nel minor tempo possibile ed esegue il task richiesto $\rightarrow$ tempo di reazione detto latenza dell’evento)
\begin{itemize}
    \item latenza di interruzione: tempo trascorso tra la ricezione dell’interruzione e l’inizio della procedura di gestione
    \item latenza di dispatch: tempo impiegato dal dispatcher per eseguire lo switch di contesto
\end{itemize}
\textbf{Scheduling Sistemi Real-Time}
\begin{itemize}
    \item \textbf{Priorità}: scheduling con priorità con prelazione visto in precedenza o varianti per garantire deadline hard e gestire task periodici
    \item \textbf{Priorità proporzionale a frequenza}: come precedente ma con priorità assegnate staticamente in base a frequenza del task periodico
    \item \textbf{Earliest Deadline First}: priorità assegnate dinamicamente (prima il task che ha scadenza più vicina, non richiede periodicità)
    \item \textbf{A quote proporzionali}: assegnata una percentuale del tempo di CPU a ogni applicazione, a ogni richiesta di utilizzo CPU si valuta se ammetterla o meno
\end{itemize}
\subsection{Scheduling in UNIX e Linux}
\subsubsection{Scheduling di UNIX Tradizionale}
\begin{itemize}
    \item Utilizza \textbf{code multiple con priorità dinamiche} e algoritmo \textbf{round-robin} per ogni coda.
    \item \textbf{Preemption} dopo 1 secondo di utilizzo continuato della CPU.
    \item Ogni processo ha una \textbf{priorità di base} che lo colloca in una banda specifica (es. Swapper, User processes) per ottimizzare l’uso dei dispositivi I/O.
    \item La priorità viene modificata in base all’utilizzo della CPU e alle preferenze dell’utente (nice), ma rimane nella stessa banda.
    \item Le priorità sono ricalcolate ogni secondo.
\end{itemize}
\subsubsection{Scheduling in O(1) di Linux}
\begin{itemize}
    \item Con \textbf{preemption}.
    \item Basato su \textbf{priorità numeriche} con due gamme:
        \begin{itemize}
            \item Priorità alta (es. 0) $\rightarrow$ quanto lungo (200 ms).
            \item Priorità bassa (es. 140) $\rightarrow$ quanto breve (10 ms).
        \end{itemize}
\end{itemize}
\textbf{Gestione dei Task}
\begin{itemize}
    \item \textbf{Task ready}: classificati come \textbf{active} (se hanno tempo residuo) o \textbf{expired} (se hanno esaurito il quanto).
    \item Il task \textbf{running} diventa \textbf{expired} quando consuma il suo quanto e non può essere rieseguito finché ci sono task active.
    \item Due array gestiscono task active e expired. La CPU è assegnata al task active con priorità massima.
    \item Gli array si scambiano i ruoli quando tutti i task \textbf{active} diventano \textbf{expired}.
\end{itemize}

\end{document}
